\chapter{Introduction} \label{chap:introduction}
\chaptersummary{Predictions are often available for many units, whereas gold-standard labels are expensive and scarce. This chapter introduces the central problem of the thesis, defines the main terms, and explains how the methodological foundations connect to the protein--protein interaction application.}

Many modern empirical studies have access to two very different kinds of information. Model predictions or other surrogate quantities can be generated cheaply for many units, while gold-standard labels are expensive and therefore available only on a smaller subset. Prediction-powered inference (PPI) addresses that setting by using a large prediction-only sample for broad coverage and a smaller labeled subset for residual correction \parencite{ppi_framework}.

The rest of the introduction keeps that setup at a high level. The aim is to make clear what is being predicted, what is being estimated, and why the protein application provides a useful test case for the broader prediction-powered workflow.

\section{Overview and Terminology}

The key idea is straightforward. A predictor can be useful for two different reasons. It may help rank observations, which is a prediction problem, and it may help estimate a scientifically meaningful population quantity when labels are scarce, which is an inference problem. PPI addresses the second task. Instead of replacing the gold-standard labels, it uses them on a smaller subset to correct what the cheap surrogate gets wrong \parencite{ppi_framework}. A tuned version of the method, written as \texttt{PPI++}, introduces a weight $\lambda$ that controls how strongly the prediction-only sample enters the corrected estimator. Informally, the weight lets the method borrow more or less aggressively from the large surrogate-labeled sample depending on how informative the surrogate is \parencite{ppi_pp}.

Throughout this thesis, the abbreviation ``PPI'' refers only to \emph{prediction-powered inference}. The biological application is written explicitly as protein--protein interaction, protein-interaction prediction, or protein--protein interaction network to avoid ambiguity. Table~\ref{tab:intro_terminology} collects the main abbreviations and working terms used early in the thesis.

\begin{table}[h!]
\centering
\small
\setlength{\tabcolsep}{6pt}
\renewcommand{\arraystretch}{1.15}
\begin{tabular}{@{}p{0.16\textwidth}p{0.74\textwidth}@{}}
\toprule
Term & Working meaning in this thesis \\
\midrule
PPI & Prediction-powered inference: use a large surrogate-labeled sample together with a smaller gold-standard labeled sample to estimate a target quantity. \\
\texttt{PPI++} & Tuned version of PPI with weight $\lambda$, used to control how much the prediction-only sample contributes after correction. \\
LLM & Large language model. In Chapter~\ref{chap:foundation_llm}, LLM judge labels are the motivating example of a cheap but noisy surrogate. \\
EIF & Efficient influence function. This is the semiparametric tool used in Chapter~\ref{chap:foundation_llm} to connect calibration ideas to efficient estimation. \\
i.i.d. & Independent and identically distributed. Chapters~\ref{chap:foundation_power} and~\ref{chap:protein_results} use this as the published planning baseline before adding dependence sensitivity. \\
Surrogate & Cheap predictive quantity available broadly, used to help estimate a target defined on gold-standard labels. \\
Pilot or calibration set & Labeled subset used either to correct the surrogate or to estimate planning quantities before a larger study is run. \\
\bottomrule
\end{tabular}
\caption[Basic terminology used in the thesis.]{Basic terminology used in the thesis.}
\label{tab:intro_terminology}
\end{table}

\section{Prediction and Inference Answer Different Questions}

One source of confusion in this area is that prediction and inference use some of the same data but answer different questions. Predictive performance concerns how well a model scores or ranks future observations. In this thesis, predictive summaries include precision-recall area under the curve (PR-AUC), receiver operating characteristic area under the curve (ROC-AUC), precision, recall, log loss, and Brier score. Inferential performance concerns whether the analysis recovers the target population quantity and quantifies its uncertainty appropriately from a limited labeled subset. In this thesis, inferential summaries include bias relative to a full-data reference, stability of corrected estimates, confidence intervals, statistical power, and required labeled sample size.

These two layers are related but not interchangeable. A predictor can have strong ranking performance and still be poorly suited for a prevalence estimate if it is badly calibrated for that inferential target. Conversely, a surrogate can be a noisy classifier and still be useful for planning if it tracks the target quantity closely on the labeled subset. When this thesis refers to statistically reliable inference, it means recovering a target population quantity and its uncertainty from a small gold-standard labeled subset plus a larger surrogate-labeled sample under the assumptions of the working model. The later chapters therefore keep predictive and inferential claims separate instead of treating every gain in one layer as a gain in the other.

\section{Thesis Components and Biological Application}

This thesis has three connected components. The first studies how to correct estimates when a cheap surrogate is available broadly but gold-standard labels are observed only on a smaller subset. The second studies how the same surrogate enters power calculations and labeled-sample planning before a full validation study is run. The third is an applied case study in protein--protein interaction prediction, where a learned surrogate score is carried into downstream estimation and planning. Chapters~\ref{chap:foundation_llm} and~\ref{chap:foundation_power} develop the two methodological foundations. Chapters~\ref{chap:protein_pipeline} and~\ref{chap:protein_results} then build and analyze the protein application.

The biological scope of that application is intentionally narrower than the phrase \emph{protein interaction} sometimes suggests. Here a positive edge means a \emph{physical protein--protein interaction} recorded in the curated benchmark resources used by the study, not a looser biological relationship such as co-expression, pathway membership, homology, or paralogy. The Intra benchmark provides predefined positive and negative edge files. The BIOGRID-human benchmark contributes curated positive physical interactions, and its negative edges are constructed as randomly sampled non-positive pairs rather than experimentally verified non-interactions. The goal is therefore not to cover every notion of biological relatedness between proteins, but to study one specific physical-interaction prediction problem under clearly defined benchmark labels \parencite{vidal2011, biogrid2019, luck2020, uniprot2025}.

This setting is a useful test case for prediction-powered methods because it has the two features that matter most for the rest of the thesis. First, physical interaction labels are biologically important and expensive to validate, so labeled data are meaningfully scarce. Second, modern sequence and annotation models can still produce large numbers of candidate interaction scores. That combination makes protein--protein interaction prediction a natural place to ask whether a strong surrogate can improve downstream estimation and planning without replacing statistical calibration.

\section[Related Work]{Related Work}

\subsection[Prediction-Powered Inference and Efficiency]{Prediction-Powered Inference and Efficient Estimation}

The statistical backbone of the thesis can be traced to earlier work on imperfect diagnostics, measurement error, missing data, and semiparametric efficiency. Those literatures asked how to recover a target parameter when the observed quantity is only a proxy for truth or when the target outcome is available only on a subset of units \parencite{rogan_gladen, rubin1976, robins1995, eif_ref, tsiatis2006}. The original PPI framework recasts that question for modern prediction systems: cheap predictions are available broadly, gold-standard labels are scarce, and the inferential goal is to combine the two without discarding the smaller labeled sample \parencite{ppi_framework}. \texttt{PPI++} makes that logic more explicit by using the tuning parameter $\lambda$ to control how strongly the prediction-only sample enters the estimator when the surrogate is informative \parencite{ppi_pp}.

The noisy-surrogate extension in \textcite{llmjudge_eif} clarifies how prediction-powered correction relates to classical misclassification correction and to efficient-influence-function arguments when the cheap signal is noisy or biased. That perspective is useful well beyond LLM evaluation. It gives the thesis a general language for treating protein-edge prediction scores as inferential surrogates rather than as stand-alone benchmark outputs.

\subsection[Power and Planning]{Prediction-Powered Planning and Power Analysis}

Most of the early prediction-powered literature is retrospective: given a predictor and a labeled subset, how should one estimate a target parameter efficiently? The planning literature asks the prospective version of the same question: before labels are collected, how does surrogate quality translate into statistical power and required labeled sample size? \textcite{ppi_power} provides the main answer used in this thesis by showing, in the one-sample i.i.d. setting, how agreement between truth and surrogate enters the asymptotic variance and the resulting labeled-sample formulas.

This design perspective is essential to the thesis because it extends the value of a surrogate beyond retrospective correction. Chapter~\ref{chap:foundation_power} treats the published i.i.d. formulas as the baseline theory. The later protein chapters then ask a narrower and more practical question: how far can those formulas be carried into a graph-structured biological setting, where dependencies between protein--protein interactions (edges) are plausible?

\subsection[Protein-Interaction Prediction]{Protein-Interaction Prediction and Benchmark Design}

The protein application draws on a separate literature on interactome resources and protein-interaction prediction. Curated resources such as BioGRID and UniProt are essential both for building datasets and for interpreting predicted interactions downstream \parencite{biogrid2019, uniprot2025}. On the predictive side, the field has moved from hand-engineered features toward deep sequence models, protein language model embeddings, and hybrid architectures that combine several information sources \parencite{sun2017, dscript2021, prottrans2022, lin2023, lee2023}.

Previous studies show that strong single-protein representations are now practically attainable. They also show that evaluation in protein-pair prediction is highly sensitive to split construction and negative-edge design \parencite{park2012, bernett2024}. This thesis follows that lesson directly. It uses the Intra benchmark with its predefined negatives, the BIOGRID-human benchmark with randomly sampled non-positive pairs, and it interprets cross-dataset comparisons cautiously because the label semantics are not identical.

Taken together, this thesis is not simply a review of semiparametric efficiency, a benchmark analysis of protein–protein interaction prediction, or a reproduction of recent methods. Rather, its contribution is to bring prediction-powered inference, prediction-powered planning, and leakage-aware protein-edge modeling into one coherent inferential framework, while clarifying both the aspects that transfer naturally and the aspects that require additional care in the protein-interaction setting.  

\section{Objectives and Improvement Criteria}

Success should be defined differently in the prediction, inference, and planning layers. Each layer has a different target, a different baseline, and a different notion of improvement. Table~\ref{tab:intro_objectives} states those distinctions explicitly.

\begin{table}[h!]
    \centering
    \small
    \setlength{\tabcolsep}{5pt}
    \renewcommand{\arraystretch}{1.15}
    \begin{tabular}{@{}p{0.12\textwidth}p{0.24\textwidth}p{0.22\textwidth}p{0.32\textwidth}@{}}
    \toprule
    Layer & Target & Baseline & Improvement criterion \\
    \midrule
    Prediction & Surrogate score for whether a protein pair interacts & Simpler models on the same leakage-aware splits & Better held-out discrimination and calibration, judged by PR-AUC, log loss, and Brier score \\
    Inference & Prevalence and related mean-type quantities on held-out edge sets & Label-only estimator and vanilla PPI & Estimates close to the full-data reference with stable uncertainty under limited labels \\
    Planning & Labeled sample size for a target effect size and power level & Label-only design and the published i.i.d. planning baseline & Fewer labels needed for the same target \\
    \bottomrule
    \end{tabular}
    \caption[Targets, baselines, and improvement criteria]{Targets, baselines, and improvement criteria for the three layers of the thesis.}
    \label{tab:intro_objectives}
    \end{table}

Two additional baseline comparisons are important for the chapters that follow. On the estimation side, the main comparison is among the classical label-only estimator, vanilla PPI, and tuned \texttt{PPI++}. On the planning side, the main comparison is between the published i.i.d. theory and the dependence-sensitive extension used in the protein application. Throughout the thesis, a claimed gain should therefore be read relative to one of these explicit baselines rather than as a generic statement that one method is simply \emph{better}.

\section{Research Questions}

The thesis is organized around four questions:
\begin{enumerate}
    \item How should the original PPI and \texttt{PPI++} framework be understood as part of a broader literature on surrogate-assisted inference, measurement error, and semiparametric efficiency?
    \item How should prediction-powered methods be interpreted when the cheap labels are themselves noisy or surrogate outcomes?
    \item How does surrogate quality enter power and sample-size planning in prediction-powered inference?
    \item In protein--protein interaction prediction, can a strong surrogate model be built and then used as the prediction layer for calibrated prediction-powered estimation and \texttt{PPI++} analysis under leakage-aware evaluation and shared-protein dependence sensitivity?
\end{enumerate}

\section{Main Contributions}

\begin{enumerate}
    \item It gives a unified interpretation of prediction-powered inference as a workflow for estimation, uncertainty quantification, and planning when labels are expensive and surrogates are cheap.
    \item It explains how the noisy-label extension in \textcite{llmjudge_eif} connects prediction-powered correction to misclassification correction, missing-data structure, and efficient influence functions.
    \item It explains how the planning formulas in \textcite{ppi_power} translate surrogate quality into labeled-sample requirements and then carries that logic into a protein-application setting.
    \item It develops a leakage-aware protein-interaction prediction workflow and shows how its fixed surrogate labels can be used for corrected estimation, labeled-budget analysis, and dependence-sensitive planning relative to explicit classical and prediction-powered baselines.
\end{enumerate}

\section{Thesis Roadmap}

Chapter~\ref{chap:foundation_llm} develops the first methodological foundation. It uses noisy LLM-as-a-judge evaluation as the motivating example for how a cheap but noisy signal can still support corrected estimation when a smaller labeled subset is available. Chapter~\ref{chap:foundation_power} develops the second foundation by turning from estimation to planning and showing how surrogate quality changes required labeled sample size before a study is run.

Chapter~\ref{chap:protein_pipeline} then turns to the protein application. It explains the biological task, the curated datasets, the leakage-aware split design, and the predictive workflow used to build the surrogate that is carried forward. Chapter~\ref{chap:protein_results} uses that frozen surrogate for three distinct purposes: predictive evaluation, corrected estimation of prevalence-like quantities, and dependence-sensitive planning. Chapter~\ref{chap:conclusion} concludes with limitations and future directions.
